% ===========================================================================
\chapter{Matrizes de Rastreabilidade da Experiência do Usuário}
\label{cap_matrizes_rastreabilidade_ux}
% ===========================================================================

Para assegurar o fechamento de ciclo metodológico entre a pesquisa qualitativa de personas e a engenharia de software da plataforma, são estabelecidas quatro matrizes de rastreabilidade bidirecional. Essas matrizes comprovam que nenhum requisito de software foi concebido sem fundamentação nas necessidades dos usuários e que todas as dores mapeadas possuem soluções técnicas implementadas.

% ===========================================================================
\section{Matriz: Dores das Personas vs. Requisitos Funcionais}
\label{sec_matriz_dores_rf}
% ===========================================================================

A Tabela \ref{tab_rastreabilidade_dores_rf} mapeia cada dor psicológica e operacional diagnosticada nos mapas de empatia em relação aos vinte e cinco requisitos funcionais (\texttt{RF01} a \texttt{RF25}) do SIGEA-GTT.

\begin{table}[!ht]
\centering
\caption{Matriz de Rastreabilidade: Dores das Personas vs. Requisitos Funcionais}
\label{tab_rastreabilidade_dores_rf}
\footnotesize
\begin{tabularx}{\textwidth}{p{2.6cm} p{4.6cm} p{1.6cm} X}
\toprule
\textbf{Persona Afetada} & \textbf{Dor Operacional / Cognitiva Mapeada} & \textbf{Requisito} & \textbf{Mecanismo de Solução no Software} \\
\midrule
\textbf{Beatriz Matos} \par (Discente) & Ansiedade aguda com cronômetro sem avisos graduais & \texttt{RF16} & Contagem regressiva visível com alertas cromáticos graduais aos 15 e 18 minutos. \\
\addlinespace
\textbf{Beatriz Matos} & Sobrecarga de leitura em prontuários físicos volumosos & \texttt{RF15} & Prontuário sintético organizado em abas estruturadas com dados em \texttt{JSONB}. \\
\addlinespace
\textbf{Beatriz Matos} & Dificuldade em memorizar os 53 gatilhos canônicos do IHI & \texttt{RF08} & Catálogo didático de gatilhos pesquisável em tempo real com filtro por módulo. \\
\addlinespace
\textbf{Beatriz Matos} & Desarticulação entre gatilhos e ferramentas de melhoria & \texttt{RF18} a \texttt{RF22} & Preenchimento integrado de Ishikawa 6M, Matriz GUT, 5W2H e ciclo PDCA. \\
\addlinespace
\textbf{Beatriz Matos} & Medo de retaliação e cultura punitiva de culpar o operador & \texttt{RF17}, \texttt{RF24} & Justificativa clínica focada em falhas sistêmicas e parecer formativo docente. \\
\midrule
\textbf{Prof. Ricardo} \par (Docente) & Exaustão física ao corrigir manualmente dezenas de papéis & \texttt{RF24} & Espelho avaliativo automatizado justapondo a entrega do aluno ao gabarito docente. \\
\addlinespace
\textbf{Prof. Ricardo} & Risco de erros em fórmulas complexas no Excel & \texttt{RF25} & Consolidação matemática automática das taxas epidemiológicas do IHI ($TI_{1000}$, $P_{EA}$). \\
\addlinespace
\textbf{Prof. Ricardo} & Turmas criadas com nomes confusos ou despadronizados & \texttt{RF10} & Máscara canônica obrigatória \texttt{"Disciplina - Turma - Período"} no cadastro. \\
\addlinespace
\textbf{Prof. Ricardo} & Fechar o semestre letivo esquecendo alunos sem notas & \texttt{RF13} & Bloqueio automático de encerramento de turma se houver atividades sem nota. \\
\midrule
\textbf{Carlos Santos} \par (Administrador) & Sobrecarga de chamados para redefinição manual de senhas & \texttt{RF02}, \texttt{RF03} & Troca obrigatória no primeiro login e autoatendimento de alteração de senha. \\
\addlinespace
\textbf{Carlos Santos} & Inconsistências de e-mail e matrículas indevidas de docentes & \texttt{RF01}, \texttt{RF05} & Validador estrito \texttt{@UfsEmail} e segregação de matrícula exclusiva para \texttt{STUDENT}. \\
\bottomrule
\end{tabularx}
\fonte{Elaborado pelo autor (2026).}
\end{table}

% ===========================================================================
\section{Matriz: Ganhos das Personas vs. Requisitos Não Funcionais}
\label{sec_matriz_ganhos_rnf}
% ===========================================================================

A Tabela \ref{tab_rastreabilidade_ganhos_rnf} relaciona os ganhos desejados pelas personas com os requisitos não funcionais (\texttt{RNF01} a \texttt{RNF10}), fundamentados nas categorias da norma ISO/IEC 25010.

\begin{table}[!ht]
\centering
\caption{Matriz de Rastreabilidade: Ganhos das Personas vs. Requisitos Não Funcionais}
\label{tab_rastreabilidade_ganhos_rnf}
\footnotesize
\begin{tabularx}{\textwidth}{p{2.6cm} p{4.6cm} p{1.6cm} X}
\toprule
\textbf{Persona} & \textbf{Ganho Almejado (Aspiração)} & \textbf{Requisito} & \textbf{Garantia Arquitetural e de Usabilidade} \\
\midrule
\textbf{Beatriz Matos} & Leitura sem fadiga visual em monitores dos laboratórios & \texttt{RNF01}, \texttt{RNF06} & Paleta hospitalar serena e leiaute homologado para resoluções HD, HD+, WUXGA e Ultrawide. \\
\addlinespace
\textbf{Beatriz Matos} & Resposta instantânea da interface sem travamentos & \texttt{RNF02}, \texttt{RNF07} & Latência de API $P_{95} \le 200\text{ ms}$, Toasts em $<100\text{ ms}$ e indicador global de carregamento. \\
\addlinespace
\textbf{Prof. Ricardo} & Prevenção contra exclusões acidentais de notas ou turmas & \texttt{RNF08} & Modais de confirmação em duas etapas com estilo visual distintivo para ações destrutivas. \\
\addlinespace
\textbf{Prof. Ricardo} & Confiabilidade matemática absoluta dos cálculos do IHI & \texttt{RNF09} & \textit{Quality Gate} com 100\% de cobertura de testes unitários e de integração (JUnit 5 / Jacoco). \\
\addlinespace
\textbf{Carlos Santos} & Conformidade rigorosa com a LGPD e dados de alunos & \texttt{RNF04} & 100\% de dados fictícios simulados e isolamento de rede contra prontuários do SUS. \\
\addlinespace
\textbf{Carlos Santos} & Visualização padronizada de usuários e turmas & \texttt{RNF05} & Paginação rígida de dez registros por página em todas as tabelas administrativas. \\
\addlinespace
\textbf{Carlos Santos} & Segurança contra acessos indevidos fora da UFS & \texttt{RNF03} & Autenticação stateless JWT com HMAC-SHA256 e senhas com hash BCrypt (custo 12). \\
\bottomrule
\end{tabularx}
\fonte{Elaborado pelo autor (2026).}
\end{table}

% ===========================================================================
\section{Matriz: Personas vs. Casos de Uso do Sistema}
\label{sec_matriz_personas_uc}
% ===========================================================================

A Tabela \ref{tab_rastreabilidade_personas_uc} apresenta o mapa de envolvimento dos três atores em relação aos vinte e dois casos de uso (\texttt{UC01} a \texttt{UC22}) do SIGEA-GTT.

\begin{table}[!ht]
\centering
\caption{Matriz de Rastreabilidade: Personas vs. Casos de Uso do Sistema}
\label{tab_rastreabilidade_personas_uc}
\footnotesize
\begin{tabularx}{\textwidth}{p{2.8cm} p{4.2cm} p{2.2cm} X}
\toprule
\textbf{Caso de Uso} & \textbf{Nome do Caso de Uso} & \textbf{Ator Primário} & \textbf{Papel e Impacto na Jornada da Persona} \\
\midrule
\texttt{UC01} a \texttt{UC03} & Autenticação e Gestão de Senha & Todos os Atores & Ingresso seguro, troca forçada inicial e autoatendimento. \\
\addlinespace
\texttt{UC04} a \texttt{UC06} & Usuários e Catálogo Didático & Administrador / Todos & Gestão de perfis e consulta didática dos 53 gatilhos. \\
\addlinespace
\texttt{UC07} a \texttt{UC10} & Gestão de Turmas e Atividades & Docente Titular & Criação canônica, matrícula em lote e trava de semestre. \\
\addlinespace
\texttt{UC11} a \texttt{UC13} & Sessão e Auditoria GTT & Discente & Rastreamento clínico sob cronômetro de 20 minutos. \\
\addlinespace
\texttt{UC14} a \texttt{UC17} & Ferramentas da Qualidade & Discente & Preenchimento de Ishikawa, GUT, 5W2H e ciclo PDCA. \\
\addlinespace
\texttt{UC18} & Submissão da Auditoria & Discente & Envio definitivo protegido por modal de confirmação. \\
\addlinespace
\texttt{UC19} a \texttt{UC22} & Avaliação e Indicadores IHI & Docente Titular & Espelho de gabarito, parecer formativo e taxas IHI. \\
\bottomrule
\end{tabularx}
\fonte{Elaborado pelo autor (2026).}
\end{table}

% ===========================================================================
\section{Matriz: Pontos de Contato da Jornada vs. Componentes Frontend}
\label{sec_matriz_touchpoints_componentes}
% ===========================================================================

A Tabela \ref{tab_touchpoints_componentes} mapeia cada ponto de contato interativo identificado nas jornadas dos usuários para os componentes \textit{Standalone} reativos desenvolvidos no Angular 18+.

\begin{table}[!ht]
\centering
\caption{Matriz de Rastreabilidade: Pontos de Contato vs. Componentes Angular 18+}
\label{tab_touchpoints_componentes}
\footnotesize
\begin{tabularx}{\textwidth}{p{2.8cm} p{2.8cm} p{4.8cm} X}
\toprule
\textbf{Ponto de Contato} & \textbf{Persona Beneficiada} & \textbf{Componente Angular 18+} & \textbf{Tecnologia de Estado / Reatividade} \\
\midrule
Tela de Acesso e Troca de Senha & Beatriz / Ricardo / Carlos & \texttt{LoginComponent} & \textit{Reactive Forms} com validadores customizados e \texttt{signal()}. \\
\addlinespace
Painel Geral do Discente & Beatriz Matos & \texttt{StudentDashboard-}\linebreak\texttt{Component} & \textit{Signals} computados (\texttt{computed()}) para cálculo de prazos. \\
\addlinespace
Área de Auditoria e Prontuário & Beatriz Matos & \texttt{AuditWorkspace-}\linebreak\texttt{Component} & \textit{Timer} RxJS com alertas de cor, abas tabuladas e modais. \\
\addlinespace
Painel de Ferramentas de Qualidade & Beatriz Matos & \texttt{QualityTools-}\linebreak\texttt{Component} & Atualização bidirecional de notas GUT com cálculo reativo. \\
\addlinespace
Gestão e Matrícula de Turmas & Prof. Dr. Ricardo Sobral & \texttt{ClassManagement-}\linebreak\texttt{Component} & Tabela Angular Material com paginação estrita e busca. \\
\addlinespace
Espelho de Correção e Gabarito & Prof. Dr. Ricardo Sobral & \texttt{EvaluationWorkspace-}\linebreak\texttt{Component} & Visão lado a lado com destaques condicionais em CSS. \\
\addlinespace
Painel de Indicadores Epidemiológicos & Prof. Dr. Ricardo Sobral & \texttt{ClassDashboard-}\linebreak\texttt{Component} & Gráficos reativos com taxas de eventos adversos do IHI. \\
\bottomrule
\end{tabularx}
\fonte{Elaborado pelo autor (2026).}
\end{table}
